\chapter{Cloud masking}\label{ch:hslcloud}
\def\tbrk{\hfil\penalty50\hfilneg}% local: a line may break before a provenance tag

\begin{onething}
A cloud classifier given HyperSL's 256-number embedding plus seven raw bands found 28.9\% of the cirrus that EnMAP's own
masks label, against 67.9\% for the same classifier given the 148 raw bands alone; HyperSL was pretrained on scenes chosen
to be nearly cloud-free, and its per-scene scaling turns the same thin cirrus into a different input in every scene.
\end{onething}

The product excludes cloud from its detections (\cref{ch:seghead}, \file{docs/lld/evaluation.md}, ``Decisions''). In
September 2026 the cloud mask came from fixed thresholds and from EnMAP's own quality layers, which are wrong over the
high plateau, and PRISMA L2D carries no cloud mask at all. A learned cloud mask on PRISMA was needed, and a foundation
model's embedding is the textbook way to build one with few labels. The experiment
\file{research/experiments/cloud-mask-hypersl/} (commits \code{fa8c1ab} and \code{b452e8d}, 2026-10-01) tried it.

\begin{analogy}
A wine taster trained only on still wines, poured at the temperature of whichever cellar they came from, and asked to
detect a faint sparkle. The sparkle is small, it was never in the training glasses, and how cold the glass is (which
changes how a sparkle feels) is set by the cellar, not the wine. Give a second taster a thermometer and the bare glass
and they do better, with no training at all on what wine is.
\end{analogy}

%======================================================================
\section{The experiment}\label{sec:hslcloud:setup}

\paragraph{Training pixels (step 1).} 97 EnMAP L2A scenes (the commit message of \code{fa8c1ab}), every spectrum
resampled onto PRISMA's band axis with a Gaussian of PRISMA's band width, because ``HyperSL's embedding shifts with band
sampling almost as much as with material'' (\cref{ch:hslnorm}). Classes 0 clear, 1 cloud, 2 cirrus, 3 snow, from EnMAP's
\code{QL_QUALITY_*} masks with three corrections agreed on 2026-10-01: scenes above 2,500\,m mean elevation contribute no
cloud or cirrus labels (EnMAP flags 51--57\% of clear plateau ground as cirrus); snow labels below 1,000\,m are dropped
(salt flats); pixels within 2 of a mask boundary are dropped. Cirrus is \code{QL_QUALITY_CIRRUS} $\ge2$; haze and shadow
are left out\tbrk\tagMr{\texttt{fa8c1ab}: step1\_training\_set.py docstring}.

\begin{sourcenote}[: the cirrus band is not there]
The classic cirrus test uses the 1.38\,$\mu$m water-vapour band, where low surfaces are dark and high ice cloud is bright.
``EnMAP L2A blanks 1343--1390\,nm \dots{} to nodata, so the 1.38\,$\mu$m cirrus band is never in the training spectra''
(the same docstring). PRISMA's 149-band axis also drops 1325--1495\,nm. Both variants below work without it, so the
comparison between them is fair, and both are harder than they would be with it.
\end{sourcenote}

\paragraph{Embeddings (step 2).} HyperSL at mask 0 on 148 bands (the 149-band axis minus 1785\,nm, which EnMAP does not
cover), each pixel scaled per band by its own scene's minimum and maximum, ``the frame HyperSL was trained in''.

\paragraph{Classifier (step 3).} A gradient-boosted tree classifier (\code{HistGradientBoostingClassifier}, 300 rounds,
classes weighted to balance, early stopping) with two feature sets:
\begin{itemize}
\item \textbf{C}, ``the chosen design'': HyperSL's 256-number embedding plus raw reflectance at 7 wavelengths (490, 560,
665, 865, 1240, 1610 and 2200\,nm);
\item \textbf{B}, ``a yardstick only'': raw reflectance on the 148 bands, no HyperSL.
\end{itemize}

\begin{definition}[Grouped cross-validation]
\emph{$k$-fold cross-validation}\index{cross-validation} splits the data into $k$ parts, trains on $k-1$ and tests on the
remaining one, $k$ times, so every sample is tested exactly once by a model that did not train on it. \emph{Grouped}
means whole groups are kept together: here an EnMAP \emph{datatake} (one overpass, cut into several tiles that share
atmosphere and sun angle) is never split between training and test.
\end{definition}
784,120 pixels from 47 datatakes, 5 folds\tbrk\tagMr{\texttt{b452e8d}: cloud-mask-hypersl/step3\_report.txt}.

%======================================================================
\section{The results, in full}\label{sec:hslcloud:results}

\begin{table}[htbp]
\centering\small
\begin{tabular}{@{}lrrrrr@{}}
\toprule
& \multicolumn{4}{c}{the classifier says} & \\
\cmidrule(lr){2-5}
EnMAP label & clear & cloud & cirrus & snow & row total \\
\midrule
clear & \textbf{341,395} & 36,215 & 4,955 & 909 & 383,474 \\
cloud & 22,708 & \textbf{244,246} & 11,452 & 2,042 & 280,448 \\
cirrus & 21,979 & 28,341 & \textbf{20,472} & 1 & 70,793 \\
snow & 1,619 & 5,025 & 36 & \textbf{42,725} & 49,405 \\
\bottomrule
\end{tabular}
\caption{Variant C, HyperSL embedding + 7 raw bands. Rows are EnMAP's labels, columns the classifier's out-of-fold
predictions\tagMr{\texttt{b452e8d}: step3\_report.txt}; row totals \tagC.}
\label{tab:hslcloud:c}
\end{table}

\begin{table}[htbp]
\centering\small
\begin{tabular}{@{}lrrrrr@{}}
\toprule
& \multicolumn{4}{c}{the classifier says} & \\
\cmidrule(lr){2-5}
EnMAP label & clear & cloud & cirrus & snow & row total \\
\midrule
clear & \textbf{342,688} & 25,158 & 12,901 & 2,727 & 383,474 \\
cloud & 27,357 & \textbf{225,595} & 25,510 & 1,986 & 280,448 \\
cirrus & 6,861 & 15,888 & \textbf{48,037} & 7 & 70,793 \\
snow & 406 & 1,135 & 16 & \textbf{47,848} & 49,405 \\
\bottomrule
\end{tabular}
\caption{Variant B, the 148 raw bands, no HyperSL; same pixels, folds and classifier\tagMr{\texttt{b452e8d}:
step3\_report.txt}.}
\label{tab:hslcloud:b}
\end{table}

\begin{definition}[Recall and precision of a class]
For a class, \emph{recall}\index{recall} is the share of pixels \emph{labelled} with it that the classifier also calls
it: the diagonal entry over its row total. \emph{Precision}\index{precision} is the share of pixels the classifier
\emph{calls} it that are labelled with it: the diagonal entry over its column total. The report prints them as ``found
$x$\% of EnMAP's'' and ``right $y$\% of the time''.
\end{definition}

\begin{worked}[label={we:hslcloud:cirrus}]{Cirrus recall and precision, from the tables}
\emph{Variant C.} The cirrus row is $21{,}979+28{,}341+20{,}472+1$:
\begin{align*}
\text{row total} &= 21{,}979+28{,}341 = 50{,}320,\\
&\quad 50{,}320+20{,}472 = 70{,}792,\quad 70{,}792+1 = 70{,}793,\\
\text{recall} &= 20{,}472/70{,}793 = 0.2892 = 28.9\%,\\
\text{column total} &= 4{,}955+11{,}452+20{,}472+36 = 36{,}915,\\
\text{precision} &= 20{,}472/36{,}915 = 0.5546 = 55.5\% .
\end{align*}
\emph{Variant B.} Same row total, 70,793.
\begin{align*}
\text{recall} &= 48{,}037/70{,}793 = 0.6786 = 67.9\%,\\
\text{column total} &= 12{,}901+25{,}510+48{,}037+16 = 86{,}464,\\
\text{precision} &= 48{,}037/86{,}464 = 0.5556 = 55.6\% .
\end{align*}
At the same precision, the raw bands find $48{,}037/20{,}472=2.35$ times as much cirrus\tagC. Of C's 50,321 missed cirrus
pixels, 21,979 were called clear and 28,341 cloud.
\end{worked}

\begin{center}\small
\begin{tabular}{@{}lrrrr@{}}
\toprule
& \multicolumn{2}{c}{recall (found)} & \multicolumn{2}{c}{precision (right)} \\
\cmidrule(lr){2-3}\cmidrule(lr){4-5}
class & C: HyperSL + 7 & B: 148 raw & C: HyperSL + 7 & B: 148 raw \\
\midrule
clear & 89.0\% & 89.4\% & 88.1\% & 90.8\% \\
cloud & \textbf{87.1\%} & 80.4\% & 77.8\% & \textbf{84.2\%} \\
cirrus & 28.9\% & \textbf{67.9\%} & 55.5\% & 55.6\% \\
snow & 86.5\% & \textbf{96.8\%} & 93.5\% & 91.0\% \\
\midrule
overall accuracy & \multicolumn{2}{c}{C 82.75\% \quad B 84.70\%\tagC} & & \\
high-plateau clear called cloud or cirrus & \multicolumn{2}{c}{C \textbf{4.34\%} \quad B 6.62\%} & & \\
\bottomrule
\end{tabular}
\end{center}

\begin{figure}[htbp]
\centering
\begin{tikzpicture}
\begin{axis}[ybar, bar width=8pt, width=0.85\linewidth, height=55mm, ymin=0, ymax=110, ylabel={recall (\%)},
  xtick={0,1,2,3}, xticklabels={clear,cloud,cirrus,snow}, enlarge x limits=0.18,
  legend style={at={(0.5,-0.18)},anchor=north,font=\small,legend columns=2}]
\addplot[fill=loss!60] table[x=idx, y=recallC] {figdata/hsl_cloud_recall.dat};
\addplot[fill=spec!70] table[x=idx, y=recallB] {figdata/hsl_cloud_recall.dat};
\legend{C: HyperSL + 7 raw bands, B: 148 raw bands}
\end{axis}
\end{tikzpicture}
\caption{Share of each EnMAP-labelled class the classifier finds, out of fold over 47 datatakes. HyperSL's features lose
most on cirrus, the thin class, and on snow.}
\label{fig:hslcloud:recall}
\end{figure}

\paragraph{Being fair.} C is not worse everywhere. It finds more of the thick cloud (87.1\% against 80.4\%), and on the
59,474 clear pixels of high-plateau scenes, where EnMAP's own masks fail, it calls fewer of them cloud or cirrus (4.34\%
against 6.62\%). Its overall accuracy is two points lower, $82.75\%$ against $84.70\%$\tagC. And the scores measure
\emph{agreement with EnMAP's masks}, not truth: the step-3 docstring warns that ``a model can only look as good as EnMAP is
right''.

\begin{sourcenote}[: one comparison the run did not make]
C has 7 raw bands and B has 148. The run did not test the 7 raw bands \emph{alone}, nor HyperSL with all 148 raw bands, so
it cannot say how much of B's advantage is ``more raw bands'' rather than ``no HyperSL''. What it does show is that adding
HyperSL's 256 numbers to 7 bands is worth far less, on cirrus, than having the other 141 bands.
\end{sourcenote}

%======================================================================
\section{Why thin cloud is hard for this model}\label{sec:hslcloud:why}

\subsection{It was never shown clouds}

The paper's pretraining data were chosen with ``cloud cover near zero'', overall below 10\% (\cref{sec:hslwhat:data}). A
reconstruction model learns what its data contain. Cloud and cirrus were deliberately kept out, so nothing in pretraining
made the embedding sensitive to them. Our own Indradhanu fine-tunes, by contrast, include cloud as a labelled class of their
segmentation head (\cref{ch:seghead}).

\subsection{The same cirrus, a different input in every scene}

Thin cirrus mostly adds a little brightness to whatever is beneath it\tagA. Call that addition $\delta$ in one band. Under
\cref{eq:hslnorm:minmax} a pixel's input changes by
\begin{align*}
x'_{\text{cirrus}}-x'_{\text{clear}} &= \frac{(x+\delta)-m}{M-m}-\frac{x-m}{M-m} && \text{the scaled value with and without }\delta\\
&= \frac{x+\delta-m-x+m}{M-m} && \text{one fraction, same denominator}\\
&= \frac{\delta}{M-m}. && x\text{ and }m\text{ cancel}
\end{align*}
The step depends on the scene's span $M-m$ in that band.

\begin{worked}[label={we:hslcloud:span}]{One cirrus layer, two scenes}
Cirrus adds $\delta=0.03$ reflectance at 865\,nm\tagA. Scene A is dark (water and forest), with $m=0.02$ and $M=0.22$, a span
of 0.20; scene B is a bright desert with $m=0.20$ and $M=0.80$, a span of 0.60\tagA.
\begin{align*}
\text{scene A: } & 0.03/0.20 = 0.15,\\
\text{scene B: } & 0.03/0.60 = 0.05 .
\end{align*}
The same cirrus is a step of 0.15 in one scene's input and 0.05 in the other's, a factor of three. With raw reflectance the
step is 0.03 in both. A classifier trained across 47 datatakes sees one physical signal at many sizes.
\end{worked}

Every band has its own span, so the cirrus signature's \emph{shape} across bands is distorted too, differently per scene
(the angle argument of \cref{we:hslnorm:angle}). And the brightest pixels of a cloudy scene are often the cloud itself: the
cloud then sets $M$, the frame, for every other pixel (\cref{we:hslnorm:anom}).

\begin{keypoint}
A cloud mask must recognise a class by an absolute property (thin ice cloud is bright, flat and high). A model that
removes absolute brightness per scene and never saw clouds has to recover both from seven raw bands bolted on afterwards.
\end{keypoint}

\paragraph{What the product does.} Cloud is a class of Indradhanu's own segmentation head, trained on labelled masks from
reflectance input; since FTH3 the release candidate excludes \emph{thick} cloud only, because its thin-cloud output called
haze cloud on set4 (\cref{ch:heavy}). The HyperSL classifier was not adopted.

\begin{drift}
The 784,120 pixels, 47 datatakes and both confusion tables belong to the 2026-10-01 run (commit \code{b452e8d}); its
training tables lived on the GPU box (\file{/wt/cloudmask/}), not in the repository. A step 4 (maps on PRISMA) was planned
in the scripts; no result of it is in the history.
\end{drift}

%======================================================================
\begin{confusions}
\confusion{HyperSL is worse at every cloud class.}{It finds more thick cloud (87.1\% against 80.4\%) and makes fewer false
calls on the high plateau. It loses on cirrus and snow, and on overall accuracy by two points.}
\confusion{The cirrus result is low because the 1.38\,$\mu$m band was missing.}{It was missing for both variants. The
difference between them is not that band.}
\confusion{Recall and precision are the same thing seen twice.}{Recall divides by the row (what EnMAP labelled), precision by
the column (what the classifier called). C and B have the same cirrus precision and very different recall.}
\confusion{These numbers say how often the mask is right.}{They say how often it agrees with EnMAP's own masks, which are
themselves wrong over the plateau.}
\end{confusions}

\begin{checkyourself}
\item\label{ex:hslcloud:1} From \cref{tab:hslcloud:c}, compute variant C's recall and precision for snow.
\item\label{ex:hslcloud:2} From \cref{tab:hslcloud:b}, compute variant B's recall and precision for cloud.
\item\label{ex:hslcloud:3} Compute both variants' overall accuracy from the tables (the diagonal over 784,120).
\item\label{ex:hslcloud:4} A third scene has span 0.40 at 865\,nm. What input step does the same cirrus ($\delta=0.03$) make?
  What span would make it 0.10?
\item\label{ex:hslcloud:5} In variant C, what share of the missed cirrus was called cloud? Why is that error less harmful to the
  product than calling it clear?
\end{checkyourself}

\nextchapter{ch:hslverdict}
